\documentclass[10pt,a4paper]{amsart}
\usepackage{amsmath,amssymb,amsthm,graphicx}
\usepackage{verbatim}
\newenvironment{reference}{\comment}{\endcomment}
%\newenvironment{reference}{}{}
\newenvironment{slogan}{\comment}{\endcomment}
\newenvironment{history}{\comment}{\endcomment}
% For commutative diagrams we use Xy-pic
\usepackage[all]{xy}

% We use 2cell for 2-commutative diagrams.
\xyoption{2cell}
% Theorem environments.
%
\theoremstyle{plain}
\newtheorem{theorem}[subsection]{Theorem}
\newtheorem{proposition}[subsection]{Proposition}
\newtheorem{lemma}[subsection]{Lemma}

\theoremstyle{definition}
\newtheorem{definition}[subsection]{Definition}
\newtheorem{example}[subsection]{Example}
\newtheorem{exercise}[subsection]{Exercise}
\newtheorem{situation}[subsection]{Situation}

\theoremstyle{remark}
\newtheorem{remark}[subsection]{Remark}
\newtheorem{remarks}[subsection]{Remarks}

\numberwithin{equation}{subsection}

% Macros
%
\def\lim{\mathop{\mathrm{lim}}\nolimits}
\def\colim{\mathop{\mathrm{colim}}\nolimits}
\def\Spec{\mathop{\mathrm{Spec}}}
\def\Hom{\mathop{\mathrm{Hom}}\nolimits}
\def\Ext{\mathop{\mathrm{Ext}}\nolimits}
\def\SheafHom{\mathop{\mathcal{H}\!\mathit{om}}\nolimits}
\def\SheafExt{\mathop{\mathcal{E}\!\mathit{xt}}\nolimits}
\def\Sch{\mathit{Sch}}
\def\Mor{\mathop{\mathrm{Mor}}\nolimits}
\def\Ob{\mathop{\mathrm{Ob}}\nolimits}
\def\Sh{\mathop{\mathit{Sh}}\nolimits}
\def\NL{\mathop{N\!L}\nolimits}
\def\CH{\mathop{\mathrm{CH}}\nolimits}
\def\proetale{{pro\text{-}\acute{e}tale}}
\def\etale{{\acute{e}tale}}
\def\QCoh{\mathit{QCoh}}
\def\Ker{\mathop{\mathrm{Ker}}}
\def\Im{\mathop{\mathrm{Im}}}
\def\Coker{\mathop{\mathrm{Coker}}}
\def\Coim{\mathop{\mathrm{Coim}}}

% Boxtimes
%
\DeclareMathSymbol{\boxtimes}{\mathbin}{AMSa}{"02}

%
% Macros for moduli stacks/spaces
%
\def\QCohstack{\mathcal{QC}\!\mathit{oh}}
\def\Cohstack{\mathcal{C}\!\mathit{oh}}
\def\Spacesstack{\mathcal{S}\!\mathit{paces}}
\def\Quotfunctor{\mathrm{Quot}}
\def\Hilbfunctor{\mathrm{Hilb}}
\def\Curvesstack{\mathcal{C}\!\mathit{urves}}
\def\Polarizedstack{\mathcal{P}\!\mathit{olarized}}
\def\Complexesstack{\mathcal{C}\!\mathit{omplexes}}
% \Pic is the operator that assigns to X its picard group, usage \Pic(X)
% \Picardstack_{X/B} denotes the Picard stack of X over B
% \Picardfunctor_{X/B} denotes the Picard functor of X over B
\def\Pic{\mathop{\mathrm{Pic}}\nolimits}
\def\Picardstack{\mathcal{P}\!\mathit{ic}}
\def\Picardfunctor{\mathrm{Pic}}
\def\Deformationcategory{\mathcal{D}\!\mathit{ef}}
\usepackage{cleveref}
\crefname{originallemma}{original lemma}{original lemmas}
\Crefname{originallemma}{Original lemma}{Original lemmas}
\makeatletter
\newlabel{lemma-dualizing-glue}{{lemma-dualizing-glue}{source}{}{sourcepointer-1}{}}
\newlabel{lemma-dualizing-glue@cref}{{[originallemma][0][]lemma-dualizing-glue}{[1][0][]source}{}{}{}}
\makeatother
\crefname{originallemma}{original lemma}{original lemmas}
\Crefname{originallemma}{Original lemma}{Original lemmas}
\makeatletter
\newlabel{lemma-dualizing-localize}{{lemma-dualizing-localize}{source}{}{sourcepointer-2}{}}
\newlabel{lemma-dualizing-localize@cref}{{[originallemma][0][]lemma-dualizing-localize}{[1][0][]source}{}{}{}}
\makeatother
\crefname{originallemma}{original lemma}{original lemmas}
\Crefname{originallemma}{Original lemma}{Original lemmas}
\makeatletter
\newlabel{lemma-dualizing-unique}{{lemma-dualizing-unique}{source}{}{sourcepointer-3}{}}
\newlabel{lemma-dualizing-unique@cref}{{[originallemma][0][]lemma-dualizing-unique}{[1][0][]source}{}{}{}}
\makeatother
\crefname{originalreference}{original reference}{original references}
\Crefname{originalreference}{Original reference}{Original references}
\makeatletter
\newlabel{more-algebra-lemma-invertible-derived}{{more-algebra-lemma-invertible-derived}{source}{}{sourcepointer-4}{}}
\newlabel{more-algebra-lemma-invertible-derived@cref}{{[originalreference][0][]more-algebra-lemma-invertible-derived}{[1][0][]source}{}{}{}}
\makeatother
\crefname{originalsection}{original section}{original sections}
\Crefname{originalsection}{Original section}{Original sections}
\makeatletter
\newlabel{section-dualizing}{{section-dualizing}{source}{}{sourcepointer-5}{}}
\newlabel{section-dualizing@cref}{{[originalsection][0][]section-dualizing}{[1][0][]source}{}{}{}}
\makeatother
\title{Gorenstein rings, invertibility of the dualizing complex and openness of the Gorenstein locus}
\author{The Stacks Project Authors}
\date{stacks-project:dualizing.tex:local-snapshot-20261012; GNU FDL 1.2 or later}
\begin{document}
\maketitle
\noindent\textit{Editorial note.} All delivered mathematical text is the author's own, copied from the pinned Stacks Project chapter source sources/stacks-dualizing/pinned-original-source.tex; no statement or proof step was rewritten, reordered or removed.\par\medskip
\noindent\textit{Editorial note.} Preamble provenance: the Stacks Project keeps one single shared preamble for every chapter; this card takes it directly from the staged sources/stacks-dualizing/preamble.tex as three verbatim spans (lines 8-12, 14-18 and 154-217), so no derived file was written and nothing was added or redefined. The wrapper class and the base packages are supplied by the assembler.\par\medskip
\noindent\textit{Editorial note.} Bibliography provenance: neither the counted statement nor its proof contains a citation (this lane rescans the whole span before writing the card), so the card carries no bibliography and no entry was reconstructed.\par\medskip
\noindent\textit{Editorial note.} Selection provenance: this card was opened under the file-based control contract (CONTROL.md), with the mechanical pre-card gate run on its statement and proof interval before the card was written.\par\medskip
\noindent\textit{Allowed standard prerequisites (editorial).}\begin{itemize}
\item The localization of a dualizing complex at a prime is a dualizing complex over the local ring. Invoked inside the counted proof (source lines 4022-4024). Transports the dualizing complex to every local ring for the local statements.
\item The uniqueness of dualizing complexes up to an invertible object. Invoked inside the counted proof (source lines 4024-4026). Turns the local Gorenstein condition into invertibility of the localized complex.
\item The locality of the property of being a dualizing complex. Invoked at the end of the counted proof (source lines 4057-4058). Lifts the local invertibility to the global statement in the second part.
\end{itemize}\medskip
\ifdefined\hypertarget\hypertarget{sourcepointer-1}{}\fi
\noindent\textit{Source reference.} Original lemma lemma-dualizing-glue: sources/stacks-dualizing/pinned-original-source.tex, source lines 2977--2977. Navigation pointer only: the counted text refers to the Stacks item lemma-dualizing-glue, which is an existing labeled result of this same chapter and is not part of this card. The pointer carries the target's own label as its number and its definition line as its span; its content is not reproduced here.\par\medskip
\ifdefined\hypertarget\hypertarget{sourcepointer-2}{}\fi
\noindent\textit{Source reference.} Original lemma lemma-dualizing-localize: sources/stacks-dualizing/pinned-original-source.tex, source lines 2950--2950. Navigation pointer only: the counted text refers to the Stacks item lemma-dualizing-localize, which is an existing labeled result of this same chapter and is not part of this card. The pointer carries the target's own label as its number and its definition line as its span; its content is not reproduced here.\par\medskip
\ifdefined\hypertarget\hypertarget{sourcepointer-3}{}\fi
\noindent\textit{Source reference.} Original lemma lemma-dualizing-unique: sources/stacks-dualizing/pinned-original-source.tex, source lines 2919--2919. Navigation pointer only: the counted text refers to the Stacks item lemma-dualizing-unique, which is an existing labeled result of this same chapter and is not part of this card. The pointer carries the target's own label as its number and its definition line as its span; its content is not reproduced here.\par\medskip
\ifdefined\hypertarget\hypertarget{sourcepointer-4}{}\fi
\noindent\textit{Source reference.} Original reference more-algebra-lemma-invertible-derived: sources/stacks-dualizing/pinned-original-source.tex, source lines 2852--2852. Navigation pointer only: the counted text refers to the Stacks item more-algebra-lemma-invertible-derived, which belongs to another chapter of the project and is not part of this card. The pointer carries the target's own label as its number because this chapter's staged snapshot does not contain the other chapter's tag table; its content is not reproduced here.\par\medskip
\ifdefined\hypertarget\hypertarget{sourcepointer-5}{}\fi
\noindent\textit{Source reference.} Original section section-dualizing: sources/stacks-dualizing/pinned-original-source.tex, source lines 2768--2768. Navigation pointer only: the counted text refers to the Stacks item section-dualizing, which is an existing labeled result of this same chapter and is not part of this card. The pointer carries the target's own label as its number and its definition line as its span; its content is not reproduced here.\par\medskip
\begin{lemma}
\label{lemma-gorenstein}
Let $A$ be a Noetherian ring.
\begin{enumerate}
\item If $A$ has a dualizing complex $\omega_A^\bullet$, then
\begin{enumerate}
\item $A$ is Gorenstein $\Leftrightarrow$ $\omega_A^\bullet$ is an invertible
object of $D(A)$,
\item $A_\mathfrak p$ is Gorenstein $\Leftrightarrow$
$(\omega_A^\bullet)_\mathfrak p$ is an invertible object of
$D(A_\mathfrak p)$,
\item $\{\mathfrak p \in \Spec(A) \mid A_\mathfrak p\text{ is Gorenstein}\}$
is an open subset.
\end{enumerate}
\item If $A$ is Gorenstein, then $A$ has a dualizing complex if and
only if $A[0]$ is a dualizing complex.
\end{enumerate}
\end{lemma}

\begin{proof}
For invertible objects of $D(A)$, see
More on Algebra, Lemma \ref{more-algebra-lemma-invertible-derived}
and the discussion in Section \ref{section-dualizing}.

\medskip\noindent
By Lemma \ref{lemma-dualizing-localize} for every
$\mathfrak p$ the complex $(\omega_A^\bullet)_\mathfrak p$ is a
dualizing complex over $A_\mathfrak p$. By definition and uniqueness
of dualizing complexes (Lemma \ref{lemma-dualizing-unique})
we see that (1)(b) holds.

\medskip\noindent
To see (1)(c) assume that $A_\mathfrak p$ is Gorenstein.
Let $n_x$ be the unique integer such that
$H^{n_{x}}((\omega_A^\bullet)_\mathfrak p)$
is nonzero and isomorphic to $A_\mathfrak p$.
Since $\omega_A^\bullet$ is in $D^b_{\textit{Coh}}(A)$
there are finitely many nonzero finite $A$-modules
$H^i(\omega_A^\bullet)$. Thus there exists some
$f \in A$, $f \not \in \mathfrak p$
such that only $H^{n_x}((\omega_A^\bullet)_f)$
is nonzero and generated by $1$ element over $A_f$.
Since dualizing complexes are faithful (by definition)
we conclude that $A_f \cong H^{n_x}((\omega_A^\bullet)_f)$.
In this way we see that $A_\mathfrak q$ is Gorenstein
for every $\mathfrak q \in D(f)$. This proves that the set
in (1)(c) is open.

\medskip\noindent
Proof of (1)(a). The implication $\Leftarrow$ follows from (1)(b).
The implication $\Rightarrow$ follows from the discussion
in the previous paragraph, where we showed that if $A_\mathfrak p$
is Gorenstein, then for some $f \in A$, $f \not \in \mathfrak p$
the complex $(\omega_A^\bullet)_f$ has only one nonzero cohomology module
which is invertible.

\medskip\noindent
If $A[0]$ is a dualizing complex then $A$ is Gorenstein by
part (1). Conversely, we see that part (1) shows that
$\omega_A^\bullet$ is locally isomorphic to a shift of $A$.
Since being a dualizing complex is local
(Lemma \ref{lemma-dualizing-glue})
the result is clear.
\end{proof}

\end{document}
